\chapter{Chapter06 — Decision Tree Hypotheses}
%%% generated by the blueprint pipeline from FABL.Chapter06.LearningAndTesting.DecisionTreeHypotheses
\section{Overview}

Declarations of the Lean module \texttt{FABL.Chapter06.LearningAndTesting.DecisionTreeHypotheses}.

\section{Declarations}

\begin{definition}[Sign and Boolean equivalence]
\lean{FABL.signBoolEquiv}
\label{FABL.signBoolEquiv}
\leanok
\uses{FABL.Sign}
Let the sign set be the multiplicative units of the integers, namely $\{1,-1\}$. The
equivalence with $\{\mathrm{false},\mathrm{true}\}$ sends $1$ to false and $-1$ to true.
\end{definition}

\begin{definition}[Decision tree hypotheses]
\lean{FABL.DecisionTreeHypothesis}
\label{FABL.DecisionTreeHypothesis}
\leanok
\uses{FABL.Sign}
For a natural number $n$, define a decision tree with coordinates in $\{0,\ldots,n-1\}$
and sign outputs recursively: a leaf stores a sign, while an internal node queries one
coordinate and has a child for each of its two possible Boolean values.
\end{definition}

\begin{definition}[Decision tree evaluation on Boolean inputs]
\lean{FABL.DecisionTreeHypothesis.eval}
\label{FABL.DecisionTreeHypothesis.eval}
\leanok
\uses{FABL.DecisionTreeHypothesis, FABL.F₂Cube, FABL.Sign}
For a decision tree on $n$ coordinates and an input $x\in(\mathbb Z/2\mathbb Z)^n$, its
value is the stored sign at a leaf. At a node querying coordinate $i$, it evaluates the
zero child if $x_i=0$ and the one child otherwise.
\end{definition}

\begin{definition}[Forgetting available coordinates]
\lean{FABL.DecisionTreeHypothesis.ofF₂DecisionTree}
\leanok
\uses{FABL.DecisionTreeHypothesis, FABL.F₂DecisionTree, FABL.Sign}
For a finite set $A$ of available coordinates and a decision tree indexed by $A$ whose
internal nodes query coordinates in $A$ and remove each queried coordinate from the
indices of both children, define an ordinary sign-valued decision tree by retaining each
leaf value and each query and recursively converting its two children.
\end{definition}

\begin{theorem}[Evaluation of a decision tree conversion]
\lean{FABL.DecisionTreeHypothesis.eval_ofF₂DecisionTree}
\leanok
\uses{FABL.DecisionTreeHypothesis.eval, FABL.DecisionTreeHypothesis.ofF₂DecisionTree, FABL.F₂Cube, FABL.F₂DecisionTree, FABL.F₂DecisionTree.eval, FABL.Sign}
\difficulty{4}
For a natural number $n$, a finite-coordinate decision tree $T$ with leaves in the sign
set, indexed by an available set of coordinates among the $n$ coordinates, and an input
$x\in\mathbb F_2^n$, the converted decision-tree hypothesis has the same value at $x$ as
$T$.
\end{theorem}

\begin{definition}[Evaluation on sign inputs]
\lean{FABL.DecisionTreeHypothesis.evaluate}
\label{FABL.DecisionTreeHypothesis.evaluate}
\leanok
\uses{FABL.BooleanFunction, FABL.DecisionTreeHypothesis, FABL.DecisionTreeHypothesis.eval, FABL.binaryCubeSignEquiv}
For a decision tree on $n$ coordinates, define its Boolean-cube function on inputs
$x\in\{ -1,1\}^n$ by converting each sign coordinate to its corresponding element of
$\mathbb Z/2\mathbb Z$ and evaluating the tree there. The function returns a sign.
\end{definition}

\begin{theorem}[Evaluation on the sign cube]
\lean{FABL.DecisionTreeHypothesis.evaluate_ofF₂DecisionTree}
\leanok
\uses{FABL.DecisionTree, FABL.DecisionTreeHypothesis.eval_ofF₂DecisionTree, FABL.DecisionTreeHypothesis.evaluate, FABL.DecisionTreeHypothesis.ofF₂DecisionTree, FABL.F₂DecisionTree.eval, FABL.Sign, FABL.binaryCubeSignEquiv}
\difficulty{1}
Let $n$ be a natural number and let $T$ be a decision tree on $n$ binary coordinates
with sign-valued leaves. For every sign vector $x\in\{-1,1\}^n$, evaluating $T$ on the
binary vector obtained by applying the inverse binary-to-sign correspondence
coordinatewise gives the same value as evaluating the associated Boolean function of $T$
at $x$.
\end{theorem}

\begin{definition}[Decision tree evaluation work]
\lean{FABL.DecisionTreeHypothesis.evaluationWork}
\label{FABL.DecisionTreeHypothesis.evaluationWork}
\leanok
\uses{FABL.DecisionTreeHypothesis}
Define the evaluation work of a sign-valued decision tree recursively: a leaf has work
$1$, and a query node has work one plus the maximum of the work of its two children.
\end{definition}

\begin{definition}[Decision tree bit encoding]
\lean{FABL.DecisionTreeHypothesis.encode}
\label{FABL.DecisionTreeHypothesis.encode}
\leanok
\uses{FABL.DecisionTreeHypothesis}
For a sign-valued decision tree, define its encoding to be the list of binary digits of
its natural-number code under the canonical encoding of the tree type.
\end{definition}

\begin{definition}[Decision tree bit decoding]
\lean{FABL.DecisionTreeHypothesis.decode}
\label{FABL.DecisionTreeHypothesis.decode}
\leanok
\uses{FABL.DecisionTreeHypothesis, FABL.SparseFourierHypothesis.natOfBits}
For a list of Boolean digits, define its decoded value to be the optional decision tree
obtained by interpreting the list as a natural number and applying the canonical partial
decoder for sign-valued decision trees on $n$ coordinates.
\end{definition}

\begin{theorem}[Decision tree encoding round trip]
\lean{FABL.DecisionTreeHypothesis.decode_encode}
\label{FABL.DecisionTreeHypothesis.decode_encode}
\leanok
\uses{FABL.DecisionTreeHypothesis, FABL.DecisionTreeHypothesis.decode, FABL.DecisionTreeHypothesis.encode}
\difficulty{1}
For every natural number $n$ and every decision tree whose leaves are signs and whose
internal nodes query coordinates in ${0,ldots,n-1}$, decoding its encoded bit list
returns that tree.
\end{theorem}

\begin{definition}[Decision tree finite representation]
\lean{FABL.DecisionTreeHypothesis.finiteRepresentation}
\label{FABL.DecisionTreeHypothesis.finiteRepresentation}
\leanok
\uses{FABL.DecisionTreeHypothesis, FABL.DecisionTreeHypothesis.decode, FABL.DecisionTreeHypothesis.decode_encode, FABL.DecisionTreeHypothesis.encode, FABL.DecisionTreeHypothesis.evaluate, FABL.DecisionTreeHypothesis.evaluationWork, FABL.FiniteHypothesisRepresentation}
For each $n$, the decision-tree hypotheses on the $n$-dimensional Boolean cube have a
finite representation consisting of a code type, a binary-string encoder and decoder
with decoding after encoding recovering every code, an evaluator returning a Boolean
function, and an evaluation-work count depending only on the code.
\end{definition}
